\section{Learning guides for moments and scores without a posterior oracle}
\label{sec:supervised-guidance}

Contract E preserves specified moments of the weighted cloud presented to the
reset. It cannot correct a cloud whose moments already differ from the desired
posterior moments. We therefore keep the reset and its analytical derivative
fixed and change the objective used to fit and select the upstream guide.
The distinction is consequential: a small error in a backward-message
regression, or small variability of importance weights, need not imply a
small error in the parameter score. The proposal below trains against the
quantities we wish to estimate, while avoiding an observed-score oracle.

The starting point is simulation. A simulated latent state is not a posterior
mean, but it is a noisy label whose conditional mean is precisely that
posterior mean. This observation gives an exact population-risk identity.
Complete-data parameter scores provide an analogous label for the observed-data
score. Neither identity asserts that a finite optimization finds its population
minimum, or that improvement averaged over simulated datasets transfers to
every particular observation sequence.

\subsection{Why one simulated trajectory can supervise conditional moments}

Let $C$ be the conditioning information available to the estimator, and let
$V=\Phi(X,C)\in\R^r$ be a vector of features of a simulated latent variable.
For filtering moments $C=(\theta,Y_{1:t})$; for a proposal informed by future
observations it may instead be $(\theta,X_{t-1},Y_{t:T})$. These are different
conditional distributions and their labels must not be interchanged. Define
$\mu(C)=\mathbb E[V\mid C]$. The estimator
$A_\eta(C,U)$ uses guide settings $\eta$ and algorithmic randomness $U$.
The features can contain state coordinates, raw products $X_jX_k$, and
selected nonlinear functions appearing in transition or observation scores.
Raw products are labels for second moments; covariance is subsequently
obtained by subtracting the outer product of the estimated mean.

\begin{assumption}[Simulation labels and information separation]
\label{ass:sg-separation}
Training pairs $(C,V)$ have the declared model law. Conditional on $C$, the
particle randomness $U$ is independent of $V$. For each fixed candidate
$\eta$, $A_\eta$ is measurable with respect to $(C,U)$ and does not receive
the latent label. Both $V$ and $A_\eta$ are square integrable. The symmetric
weight matrix $W(C)$ is positive semidefinite and specified without access to
the latent evaluation label beyond $C$; it is the same for every candidate.
All displayed weighted second moments are finite. Learned weights and settings are treated
as fixed when evaluating a fresh validation or test dataset.
\end{assumption}

\begin{proposition}[Moment supervision without known posterior moments]
\label{prop:sg-moment-risk}
Under Assumption~\ref{ass:sg-separation}, with
$\|v\|_W^2=v^\top Wv$,
\begin{equation}
\mathbb E\|A_\eta-V\|_W^2
=\mathbb E\|A_\eta-\mu(C)\|_W^2
 +\mathbb E\|V-\mu(C)\|_W^2.
\label{eq:sg-moment-risk}
\end{equation}
Consequently the supervised population loss and the actual conditional-moment
MSE have the same minimizers over any fixed candidate family.
\end{proposition}
\begin{proof}
Write $A_\eta-V=(A_\eta-\mu)-(V-\mu)$ and expand the square.
The mixed term has conditional expectation
\[
\mathbb E[(A_\eta-\mu)^\top W(V-\mu)\mid C,U]
=(A_\eta-\mu)^\top W\,
  \mathbb E[V-\mu\mid C,U]=0.
\]
The last equality uses conditional independence and the definition of $\mu$.
Taking expectations gives the identity. Its second term is independent of
$\eta$, proving the minimizer statement.
\end{proof}

For example, if $X\mid C$ is scalar normal with mean $m$ and variance $v$,
then $\mathbb E[(a-X)^2\mid C]=(a-m)^2+v$.
Replacing $X$ by $X^2$ gives a label for $m^2+v$; its irreducible conditional
variance is $2v^2+4m^2v$. Higher-order features therefore bring additional
label noise. We choose them because the score uses them, rather than assuming
that matching every available moment must help. A moment-matching objective
alone gives no theorem about the derivative of the finite LEDH program.

\subsection{Complete-data labels target the actual LEDH score risk}

Fix a parameter $\theta$ and a joint density $p_\theta(x,y)$ with a
parameter-independent dominating measure. Suppose differentiation can pass
under the integral defining $p_\theta(y)$, the relevant scores are square
integrable, and $p_\theta(y)>0$. Define
\begin{equation}
 H_\theta(x,y)=\nabla_\theta\log p_\theta(x,y),\qquad
 s_\theta(y)=\nabla_\theta\log p_\theta(y).
\label{eq:sg-score-definitions}
\end{equation}
The latent path is held fixed in $H_\theta$. Differentiating the simulated
path itself with common base noise generally computes a different quantity.

\begin{proposition}[Score supervision and paired comparisons]
\label{prop:sg-score-risk}
Under the preceding regularity conditions,
$\mathbb E[H_\theta(X,Y)\mid Y]=s_\theta(Y)$. For any estimator
$\widehat s_{\theta,\eta}(Y,U)$ satisfying the separation and integrability
conditions of Assumption~\ref{ass:sg-separation},
\begin{equation}
 \mathbb E\|\widehat s_{\theta,\eta}-H_\theta\|_W^2
 =\mathbb E\|\widehat s_{\theta,\eta}-s_\theta(Y)\|_W^2+C_\theta,
 \quad C_\theta=\mathbb E\|H_\theta-s_\theta(Y)\|_W^2.
\label{eq:sg-score-risk}
\end{equation}
For two fixed candidates $a,b$, the expectation of their difference in
supervised losses equals their difference in true score MSE.
\end{proposition}
\begin{proof}
For almost every $y$ with positive marginal density,
\[
 \int H_\theta(x,y)p_\theta(x\mid y)\,\dd x
 =\frac{1}{p_\theta(y)}\int\nabla_\theta p_\theta(x,y)\,\dd x
 =\frac{\nabla_\theta p_\theta(y)}{p_\theta(y)}.
\]
This is the stated conditional identity. Apply
Proposition~\ref{prop:sg-moment-risk} with $C=(\theta,Y)$ and $V=H_\theta$.
Subtracting the identities for $a$ and $b$ cancels $C_\theta$.
\end{proof}

Here $\widehat s$ can be the analytical total derivative of the same frozen
finite LEDH likelihood program used for evaluation. The result does not
replace that derivative by a smoothing estimator: complete-data scores are
training labels only. Nor does it require $\widehat s$ to be unbiased.
The risk includes both squared bias and particle variance. If $\theta$ is
drawn from a specified training distribution, integrating the identity over
that distribution gives the corresponding average-risk statement. A fit at
one $\theta$ establishes no coverage of other parameter regions.

The population identity does not make an in-sample loss difference an unbiased
assessment of a selected candidate. Select settings using calibration and
validation data, freeze them, and assess them on independent test sequences.
For each sequence share its label and, where possible, particle seeds between
candidates. Average particle replicates within a sequence before computing
uncertainty across independent sequences. Reusing a seed to generate both the
latent label and the particle approximation can invalidate the proof.

\subsection{Exact labels for the SIR scope actually being computed}

The active continuous SIR comparison has
$X_0\sim\mathcal N(m_0,I)$,
$X_t\mid X_{t-1}\sim\mathcal N(f_\theta(X_{t-1}),I)$ and
$Y_t\mid X_t\sim\mathcal N(BX_t,\sigma_\theta^2 I_9)$,
where $B$ extracts infectious coordinates and
$\sigma_\theta=10e^{\theta_3}$. The first two parameters are log scales of
infection and removal coefficients in the existing RK4 transition.
This is the continuous law used by the shared consumer, not an assertion
that a simulator which clips susceptible states has the same joint density.
The initial distribution is parameter independent in this scope.

\begin{proposition}[Fixed-path SIR score labels]
\label{prop:sg-sir-label}
Let $r_t=X_t-f_\theta(X_{t-1})$ and $e_t=Y_t-BX_t$.
The complete-data label in these three parameter coordinates is
\begin{equation}
 H_j=\sum_{t=1}^T
       r_t^\top\partial_{\theta_j}f_\theta(X_{t-1}),\quad j=1,2,
 \qquad
 H_3=\sum_{t=1}^T(\|e_t\|^2/\sigma_\theta^2-9).
\label{eq:sg-sir-label}
\end{equation}
\end{proposition}
\begin{proof}
The transition log density contributes $-\tfrac12 r_t^\top r_t$ and a
parameter-independent normalizer. At fixed adjacent states,
$\partial_{\theta_j}r_t=-\partial_{\theta_j}f_\theta(X_{t-1})$, giving
the first expression. The observation term is
$-9\log\sigma_\theta-\|e_t\|^2/(2\sigma_\theta^2)$ plus a constant.
Since $\partial_{\theta_3}\log\sigma_\theta=1$ and $BX_t$ is fixed,
its derivative is the second expression. Sum over time; the initial law
contributes zero.
\end{proof}

The implementation must reuse the consumer's transition and its analytical
parameter tangent with zero incoming state tangent. This prevents a second,
inconsistent SIR integrator from defining the labels. A separate fixed-path
finite-difference check verifies the sum of joint log densities. For an
alternative clipped model, one must instead use a correct mixed-measure
density or a continuous pre-clipping augmentation with the correct observation
and transition laws. Simply evaluating a Gaussian density on clipped states
would violate the assumptions of Proposition~\ref{prop:sg-score-risk}.

\subsection{Changing controlled SMC and KDM at their actual consumers}

Controlled SMC supplies a useful initial policy through the backward
least-squares recursion of \citet[Section 4.2, equations (20)--(21)]{heng2020controlled}.
We retain that recursion and evaluate settings such as basis, iteration count
and guide strength by the moments or analytical scores produced downstream.
The first implementation can select a finite, declared bank of such settings;
continuous optimization of every control coefficient is a further method,
requiring its own derivative and optimization checks. A smaller backward
residual is an explanatory measurement, not the supervised objective.

For KDM, \citet[Section 5, equation (18)]{younis2023} train density NLL and
also describe subsequent refinement of the weighted mean by MSE; Appendix
D.1, equation (23), defines that mean loss. Their released full-sequence
trainer evaluates the chosen loss on latent-state labels. Replacing the mean
feature by selected features gives Proposition~\ref{prop:sg-moment-risk};
using the actual LEDH score and complete-data labels gives
Proposition~\ref{prop:sg-score-risk}. These are adaptations. Our positive-RBF
SIR message export is KDM-inspired and does not become original neural KDM
training merely by changing its selection objective.

The consumer matters as much as the loss. If a Gaussian base is multiplied
by a Gaussian-mixture message, each product component has mean $m_k$ and
covariance $C_k$, and normalized product weight $\alpha_k$. Its moments are
\begin{equation}
 \bar m=\sum_k\alpha_km_k,\qquad
 \bar C=\sum_k\alpha_k\{C_k+(m_k-\bar m)(m_k-\bar m)^\top\}.
\label{eq:sg-product-moments}
\end{equation}
This follows by expanding
$\mathbb E[(X-\bar m)(X-\bar m)^\top\mid k]$ and averaging over $k$.
Averaging component precision matrices omits the displayed between-component
dispersion. We can use $\mathcal N(\bar m,\bar C)$ as the proposal while
retaining all analytical derivatives of these moments. It is a moment-matched
Gaussian, not a draw from the full product mixture. Consequently the
importance denominator must be this Gaussian density, followed by the
existing flow Jacobian correction. The original target density is unchanged.

\begin{proposition}[Admissible quadratic guidance]
\label{prop:sg-quadratic-domain}
Multiplying $\mathcal N(m,Q)$, $Q\succ0$, by
$\exp(-x^\top\Lambda x/2+b^\top x)$ with symmetric $\Lambda$
defines a normalizable Gaussian density if and only if
$Q^{-1}+\Lambda\succ0$.
\end{proposition}
\begin{proof}
The exponent has quadratic precision $P=Q^{-1}+\Lambda$ and linear term
$Q^{-1}m+b$. When $P\succ0$, completing the square gives covariance
$P^{-1}$ and mean $P^{-1}(Q^{-1}m+b)$ and a finite Gaussian integral.
If $P$ has a negative eigenvalue, the integral grows exponentially along
that eigendirection. If it is positive semidefinite and singular, integration
along a null direction diverges whether the corresponding linear coefficient
is zero or nonzero. These exhaust the symmetric cases.
\end{proof}

Thus $\Lambda\succeq0$ is sufficient but unnecessarily restrictive. The
optional unprojected quadratic candidate must pass the actual total-precision
check at every consumed covariance. Failure rejects the candidate; it does
not authorize clipping or silently changing the target. Product-moment and
unprojected-quadratic settings are evaluated separately, so that changed
representation and changed fitting objective can be distinguished.

\subsection{Executable training and evaluation order}

\begin{enumerate}[leftmargin=*]
\item Declare the exact model law, parameter coordinates and training region;
choose moment or score loss and its frozen scaling. Partition simulation seeds
into calibration, validation and test sets before fitting.
\item Simulate $(X,Y)$ and compute fixed-path labels. Pass only $Y$, the
declared parameter and independent fit randomness to cSMC/KDM. Latent paths
are retained solely in the label calculation.
\item Fit backward messages on observations; construct each declared guide
candidate. For each candidate run the shared full analytical LEDH executor
with independent particle noise, Contract E and OT $(1,5)$. Include the
unchanged and unguided comparators. Compute losses from the actual consumed
proposal and actual final score, rather than a detached guide-only prediction.
\item Select guide settings on calibration/validation data, with numerical
validity checked before loss. Freeze all choices and evaluate independent
test sequences. Keep latent labels and particle randomness independent
conditionally on the observations and parameter throughout.
\item Report paired score-loss differences and their uncertainty, separately
from moment error and message fit. A finite bank implements supervised
selection; it does not establish globally optimal controls, neural capacity
training, or canonical tuning admission.
\end{enumerate}

Offline fitting parameters remain fixed in the analytical score recursion.
If a future implementation makes guide coefficients explicit functions of
$\theta$ inside the value program, their derivatives must be included; a
frozen-guide partial derivative must not be called that program's total
derivative. Optimization over guide settings does not require differentiating
the score: a finite candidate bank already supplies a valid first experiment.
For gradient optimization, the required mixed derivatives concern training
parameters and must be distinguished from the analytical model-score path.

\subsection{The first SIR implementation and its limits}
\label{sec:sg-pilot}

The first implementation selects a finite bank of guide settings using the
analytical score produced by the shared LEDH executor. It includes the UKF
covariance lifecycle, flow, Contract E reset and both diagonal and pairwise
GenUT corrections. The kernel alternative uses the fitted positive-RBF
message described above. It does not train a new neural KDM. Direct moment
supervision is implemented as a loss, but the executed bank selection uses
score loss; terminal first- and second-moment errors are explanatory only.

The short SIR experiments use 64 particles at $\theta=(0,0,0)$, two
calibration sequences, two validation sequences and four untouched test
sequences, each evaluated with two independent particle streams. OT remains
$(1,5)$, with 32 terminal balancing updates also held fixed. Computation uses
double precision on the GPU as a diagnostic
exception. These small experiments test the learning signal and the actual
consumer, rather than establish settings for a longer horizon or a different
parameter region.

At $T=2$, the unprojected quadratic guide was initially rejected by an
implementation check requiring $\Lambda\succeq0$. Proposition
\ref{prop:sg-quadratic-domain} shows why that condition is unnecessarily
restrictive. An explicit optional setting now checks positive total precision
instead, while preserving the old default. Tests cover both an admissible
indefinite guide and a nonnormalizable proposal. The failed rows are retained
as implementation failures; they are not evidence against unprojected guidance.

After that repair, a fresh $T=5$ experiment completed all 96 candidate runs
with valid stage diagnostics. Each of its six candidates also passed a
directional finite-difference check of the frozen finite program. Table
\ref{tab:sg-pilot} reports the scaled supervised loss
$\sum_{j=1}^3((\widehat s_j-H_j)/a_j)^2$, where
$a_j=\max\{1,\sqrt{\operatorname{mean}_{\rm calibration}H_j^2}\}$
is fixed before validation. Thus these entries include irreducible label
noise; they are not measured squared errors against a known observed score.

\begin{table}[htbp]
\centering
\begin{tabular}{lrr}
\toprule
Guide & Validation mean loss & Test mean loss \\
\midrule
None & 1653.53 & 65902.87 \\
cSMC, half strength & 4515.01 & 23443.86 \\
cSMC, projected & 3136.57 & 9575.05 \\
cSMC, unprojected & 3151.98 & 9454.93 \\
Kernel, positive information & 3328.31 & 15098.74 \\
Kernel, product moments & 3592.32 & 10847.56 \\
\bottomrule
\end{tabular}
\caption{Descriptive $T=5$, $N=64$ supervised score losses. Particle
replicates are averaged within each sequence, then independent sequences
are averaged. The validation and test sets contain two and four sequences,
respectively. No statistical ranking follows from this table.}
\label{tab:sg-pilot}
\end{table}

Unguided LEDH has the lowest validation loss. Selection within each guide
family nominates projected cSMC and positive-information kernel guidance,
but neither nomination justifies replacing the unguided method. All guided
test means are lower than the unguided test mean, yet every guided candidate
has higher loss on one of the four test sequences. Most of the mean difference
comes from the observation-noise score and a single sequence on which the
unguided loss is very large. The separate validation and test comparisons
therefore do not establish a reliable ranking. They also do not invalidate
the population identity: a small, noisy sample need not order candidates by
their population risk.

The next scientific comparison needs fresh simulations, more independent
sequences and its own scope-specific calibration of the remaining LEDH and
guide controls. The same fixed OT choice can be retained. More particle
replicates alone would not resolve uncertainty across datasets. The
implementation plan and full paired results are preserved in
\path{docs/plans/ledh-supervised-guidance-20261019.md} and
\path{docs/plans/artifacts/ledh-supervised-guidance-20261019/result.md}.
